% Generated draft: requires mathematical and source review.
\documentclass[11pt]{article}
\usepackage[margin=25mm]{geometry}
\usepackage{fontspec}
\setmainfont{Latin Modern Roman}
\usepackage{amsmath,amssymb,mathtools,unicode-math}
\setmathfont{Latin Modern Math}
\usepackage{enumitem,needspace,xurl,hyperref}
\setlength{\parindent}{0pt}
\setlength{\parskip}{5pt}
\setlength{\emergencystretch}{3em}
\newcommand{\N}{\mathbb{N}}
\newcommand{\Z}{\mathbb{Z}}
\newcommand{\Q}{\mathbb{Q}}
\newcommand{\R}{\mathbb{R}}
\newcommand{\C}{\mathbb{C}}
\newcommand{\F}{\mathbb{F}}
\newcommand{\A}{\mathbb{A}}
\newcommand{\PP}{\mathbb{P}}
\newcommand{\E}{\mathbb{E}}
\newcommand{\eps}{\varepsilon}
\newcommand{\dd}{\,\mathrm{d}}
\begin{document}
\section*{UnsolvedMath ID: 30003610}
\textbf{Source problem code:} catalogue entry
\subsection{Definitions and mathematical statement}
The supplied question does not specify a ground field. Conditional on choosing a field \(k\), let \(X\) be a smooth projective \(k\)-variety and let \(\operatorname{Mor}_k(\mathbb P^1_k,X)\) denote the scheme of morphisms from \(\mathbb P^1_k\) to \(X\). Its functor of points sends a \(k\)-scheme \(T\) to
\[
\operatorname{Hom}_T(\mathbb P^1_k\times_k T, X\times_k T).
\]
For a scheme \(Y\), write \(\operatorname{Irr}(Y)\) for the set of irreducible components of its underlying topological space.

The supplied problem asks: for a smooth projective variety \(X\), determine the number
\[
\#\operatorname{Irr}\bigl(\operatorname{Mor}_k(\mathbb P^1_k,X)\bigr).
\]
Here the entire morphism scheme is written: no degree or curve class is fixed, constant morphisms are included, and morphisms are not identified under automorphisms of \(\mathbb P^1_k\). These conventions explain the displayed notation; the source record does not establish whether a more restricted space was intended.

The ground field, its characteristic, and whether components are to be counted over \(k\) or after extension to an algebraic closure remain unspecified. If geometric components are intended, the expression to count is instead
\[
\#\operatorname{Irr}\bigl(\operatorname{Mor}_{\overline{k}}(\mathbb P^1_{\overline{k}},X_{\overline{k}})\bigr).
\]
No finiteness assertion is supplied. The accompanying literature assessment concerns Fano varieties, whereas the stated question concerns all smooth projective varieties; a Fano hypothesis cannot be inferred from the statement. These unresolved qualifications prevent a uniquely specified counting problem.
\subsection{Short English statement}
Given a smooth projective variety \(X\), determine how many irreducible components its space of maps from \(\mathbb P^1\) has. The supplied statement leaves the ground field and component convention unspecified, and does not restrict the degree or impose a Fano hypothesis.
\subsection{Sources}
\begin{itemize}
\item Local registry: \texttt{data/unsolvedmath/problems.json}, supplied record with permanent UnsolvedMath ID \(30003610\), source code \texttt{catalogue entry}, titled ``Components of Rational-Curve Mapping Spaces.'' The record gives proposed year 2017 and names no individual proposer.
\item Supplied report citation: \emph{Algebraic Geometry: Birational Classification, Derived Categories, and Moduli Spaces} (2018), Oberwolfach Reports, pp. 2703--2767, DOI \url{https://doi.org/10.4171/owr/2017/45}. The literature-source label is ``Oberwolfach Report 45 2017.'' The supplied year labels are retained without reconciliation.
\item Supplied literature source: ``Lehmann Tanimoto 2019,'' \url{https://doi.org/10.1112/S0010437X19007103}. No article title or more specific supporting passage is supplied.
\item Literature checked at: \texttt{2026-08-21}. This is supplied dated curation, not an independently verified assessment. The record labels the problem open and states that no general formula is available for arbitrary Fano varieties, while known results cover important families and bounds. Its background explicitly describes this as dated literature triage and requests independent verification of time-sensitive status and cited claims.
\end{itemize}
\end{document}
